% !TEX TS-program = xelatex
% !TEX encoding = UTF-8 Unicode
% !Mode:: "TeX:UTF-8"

\documentclass{resume_en_tm}
% \usepackage{zh_CN-Adobefonts_external}
% \usepackage{zh_CN-Adobefonts_internal}
\usepackage{linespacing_fix}
\usepackage{cite}

\begin{document}
\pagenumbering{gobble}

\name{Minghan Wang}

\contactInfo{wangmh1994@gmail.com}{PhD Candidate, Monash University}{GitHub: yuriak}

\section{Working Experience}
\datedsubsection{\textbf{Huawei 2012 Lab --- Text Machine Translation Lab}\ \ \ \  \textit{Algorithm Engineer}}{2020.1--2023.6}
Conducted research and engineering development for Huawei Translate, Huawei's machine translation product, focusing on real-time speech translation, OCR-based multimodal translation, and translation quality enhancement.
\begin{itemize}
  \item Developed streaming ASR and speech translation components for low-latency real-time subtitle translation.
  \item Improved OCR translation algorithms for production-scale multimodal translation scenarios.
  \item Built automatic post-editing modules to enhance translation quality and support localization workflows.
  \item Participated in WMT and IWSLT shared tasks from 2020 to 2022, achieving first- and second-place results.
  \item Led research initiatives on non-autoregressive translation and speech translation, bridging academic methods with product deployment.
\end{itemize}

\datedsubsection{\textbf{Huawei 2012 Lab --- RAMS Lab}\ \ \ \ \textit{Algorithm Engineer}}{2019.1--2020.1}
Developed enterprise-level intelligent customer service systems based on retrieval and knowledge-graph technologies, improving internal knowledge access and cross-departmental question answering.
\begin{itemize}
  \item Built intent recognition and slot filling models to support robust query understanding in low-resource scenarios.
  \item Integrated retrieval-based and knowledge-graph-based QA pipelines into a unified intelligent service framework.
  \item Improved the accuracy, latency, and maintainability of production NLP and information retrieval modules.
\end{itemize}


\section{Education Experience}
\datedsubsection{\textbf{Monash University}\\ Doctor of Philosophy in Computer Science}{2023.7--Present}
\begin{itemize}
    \item Thesis: Inference as a Temporal Resource in Large Language Models \\(supervised by \textbf{Gholamreza Haffari, Ehsan Shareghi, Trang Vu})
    \item Thesis overview: Developing temporally adaptive and inference-efficient LLMs for time-sensitive tasks, covering temporal awareness, streaming processing, and latent-space reasoning.
    \item Teaching experience: Teaching Assistant for FIT5217 Natural Language Processing.
\end{itemize}

\datedsubsection{\textbf{University of Melbourne}\\ Master of Science in Computer Science}{2017.1--2019.1}
\begin{itemize}
    \item GPA: 3.75/4.0, \textbf{Top 15\% ranking, with distinction}
    \item Thesis: Learning Algorithms in Quantitative Trading (supervised by \textbf{Trevor Cohn})
    \item Explored reinforcement learning and NLP-based market modeling, developing policy-gradient trading agents that combine quantitative signals with textual financial news.
\end{itemize}

% \datedsubsection{\textbf{Dongbei University of Finance and Economics}\\ Bachelor of Management in Information Management and Information Systems}{2012.9--2016.7}
% \begin{itemize}
%     \item GPA: 3.0/4.0, \textbf{Top 20\% ranking}
%     \item Thesis: Machine-Learning-Based Online Evaluation System for Classroom Quiz Assessment \\(supervised by \textbf{Hong Li})
%     \item Conducted research on NLP-based social network analysis, including user profiling, social relation graph analysis, and public opinion monitoring.
% \end{itemize}


\section{Publications}
\begin{itemize}[parsep=0.2ex]
    \item GTS: Inference-Time Scaling of Latent Reasoning with a Learnable Gaussian Thought Sampler (EMNLP 2026)
    \item Towards Inference-time Scaling for Continuous Space Reasoning (ACL 2026)
    \item Discrete Minds in a Continuous World: Do Language Models Know Time Passes? (EMNLP 2025)
    \item Proverbs Run in Pairs: Evaluating Proverb Translation Capability of Large Language Model (ACL 2025)
    \item SpeechDialogueFactory: Generating High-Quality Speech Dialogue Data to Accelerate Your Speech-LLM Development (Interspeech 2025)
    \item Exploring the Potential of Multimodal LLM with Knowledge-Intensive Multimodal ASR (EMNLP 2024)
    \item Conversational SimulMT: Efficient Simultaneous Translation with Large Language Models (IWSLT 2025)
    \item Simultaneous Machine Translation with Large Language Models (ALTA 2024)
\end{itemize}

\end{document}